\chapter{Contexte général et étude de l'existant}

\section*{Introduction}
\addcontentsline{toc}{section}{Introduction}
Ce premier chapitre situe le projet dans son contexte métier. Il présente la problématique de la gestion d'un restaurant, analyse les solutions existantes sur le marché, puis dégage les objectifs et le périmètre fonctionnel retenus pour Tastify. Ces éléments servent de fondation à la conception détaillée présentée au chapitre suivant.

\section{Problématique}
Dans un restaurant, plusieurs acteurs interviennent simultanément : le gérant, les serveurs, les cuisiniers et les clients. Chacun manipule des informations différentes, mais toutes ces informations sont liées. Une commande prise en salle concerne une table, des plats, des ingrédients, un paiement et parfois un avis client.

Lorsque ces informations sont gérées séparément, le restaurant perd en efficacité. Une rupture de stock peut être découverte trop tard, une commande peut être mal transmise, une réservation peut entrer en conflit avec une table déjà occupée, et les retours clients peuvent rester inexploités. La problématique principale est donc la suivante :

\begin{quote}
Comment concevoir une solution web centralisée, accessible et maintenable permettant de gérer les opérations d'un restaurant tout en améliorant la coordination entre les acteurs ?
\end{quote}

\section{Etude de l'existant}
Plusieurs solutions commerciales existent pour la gestion de restaurant. Parmi les plus connues, on retrouve Lightspeed Restaurant, Toast POS, Revel Systems et L'Addition. Ces outils proposent des fonctionnalités avancées de caisse, de commande ou de reporting. Cependant, ils présentent souvent des limites pour les petites et moyennes structures : coût mensuel élevé, dépendance à du matériel spécifique, disponibilité limitée selon les pays, ou faible adaptation aux besoins locaux.

Tastify se positionne comme une solution web modulaire, open-source dans son approche technique, et conçue pour couvrir les besoins principaux sans imposer d'infrastructure lourde.

\begin{table}[H]
\centering
\caption{Comparaison synthétique des solutions existantes}
\begin{tabular}{p{3cm}p{4cm}p{4cm}p{3cm}}
\toprule
\textbf{Solution} & \textbf{Points forts} & \textbf{Limites} & \textbf{Positionnement} \\
\midrule
Lightspeed & Caisse, stocks, rapports & Coût élevé, dépendance commerciale & Solution cloud complète \\
Toast POS & Commandes et cuisine & Peu adapté au marché local & Solution orientée POS \\
Revel Systems & ERP riche & Complexité et coût & Solution entreprise \\
L'Addition & Caisse restaurant & Couverture fonctionnelle variable & Solution spécialisée \\
Tastify & Modules intégrés, web, temps réel & Projet académique évolutif & Solution PFA adaptée PME \\
\bottomrule
\end{tabular}
\end{table}

À l'issue de cette analyse, deux constats orientent le projet. D'une part, aucune des solutions étudiées ne combine, à un coût accessible, la gestion intégrée salle-cuisine, la réservation en ligne et l'exploitation des avis clients. D'autre part, l'adaptation au contexte local (devise, langue, simplicité de déploiement) reste limitée. Tastify cherche à répondre à ces deux points.

\section{Objectifs du projet}
Les objectifs principaux de Tastify sont les suivants :

\begin{itemize}[leftmargin=1.2cm]
  \item centraliser la gestion du menu, des tables, des commandes, des stocks et des employés ;
  \item améliorer la coordination entre les serveurs et la cuisine grâce au temps réel ;
  \item proposer un portail client pour consulter le menu, réserver, payer et suivre la fidélité ;
  \item analyser les avis clients afin d'aider le gérant à suivre la satisfaction ;
  \item construire une architecture maintenable, testable et déployable avec Docker.
\end{itemize}

\section{Périmètre fonctionnel}
Le périmètre retenu couvre quatre espaces d'usage :

\begin{itemize}[leftmargin=1.2cm]
  \item le back-office gérant : menu, stocks, ressources humaines, configuration et tableaux de bord ;
  \item l'espace serveur : plan de salle, prise de commande, suivi des tables et encaissement ;
  \item l'espace cuisine : Kitchen Display System et suivi de préparation des plats ;
  \item le portail client : menu public, réservation, compte client, fidélité, paiement et avis.
\end{itemize}

Dans l'implémentation actuelle, ces quatre espaces sont organisés autour de \textbf{deux applications React} : une application \emph{back-office} (staff) qui regroupe les rôles gérant, serveur et cuisinier, et une application \emph{client} dédiée au portail public. Ce choix réduit la duplication de code et simplifie le déploiement, tout en conservant une séparation claire des accès par rôle.

\section{Gestion des risques du projet}
Le développement d'un système couvrant autant de modules présente des risques techniques et organisationnels qu'il convient d'anticiper. Le tableau suivant recense les principaux risques identifiés ainsi que leur plan d'atténuation.

\begin{table}[H]
\centering
\small
\caption{Matrice des risques du projet}
\begin{tabular}{p{2.6cm}p{4.2cm}p{1.6cm}p{4.6cm}}
\toprule
\textbf{Catégorie} & \textbf{Description} & \textbf{Niveau} & \textbf{Atténuation} \\
\midrule
Technique & Latence ou coupures réseau dégradant la synchronisation WebSocket du KDS. & Moyen & Reconnexion automatique côté client et affichage du dernier état connu. \\
Adoption & Réticence du personnel à abandonner les procédures manuelles. & Moyen & Interfaces simples (parcours en peu de clics) et accompagnement à la prise en main. \\
Intelligence applicative & Historique de données insuffisant au lancement pour des recommandations pertinentes. & Faible & Démarrage par règles simples (catégorie, disponibilité), enrichissement progressif. \\
Organisationnel & Complexité cumulée des modules pouvant entraîner un dépassement du calendrier. & Élevé & Périmètre MVP priorisé ; modules avancés traités en second. \\
Sécurité & Accès non autorisé entre rôles. & Moyen & Authentification JWT, permissions par rôle et tests dédiés. \\
\bottomrule
\end{tabular}
\end{table}

\section{Méthodologie et planification}
Le projet a été conduit de manière itérative, chaque itération livrant un périmètre fonctionnel testable. Cette approche a permis de garder le contrôle sur un système multi-modules sans perdre la cohérence d'ensemble.

\begin{table}[H]
\centering
\caption{Phases de réalisation du projet}
\begin{tabular}{p{2cm}p{9cm}p{2.5cm}}
\toprule
\textbf{Phase} & \textbf{Contenu} & \textbf{Durée} \\
\midrule
Phase 1 & Fondations : environnement, modèles Django, API REST de base, authentification JWT et rôles. & 2 semaines \\
Phase 2 & Back-office gérant : menu, stocks, employés et tableaux de bord. & 2 semaines \\
Phase 3 & Salle et cuisine : prise de commande, plan de salle, KDS temps réel via WebSocket. & 2 semaines \\
Phase 4 & Portail client : menu en ligne, réservation, paiement et partage de l'addition. & 2 semaines \\
Phase 5 & Fidélité et avis : points, paliers, récompenses et analyse de sentiment. & 2 semaines \\
Phase 6 & Tests, intégration, conteneurisation Docker et finalisation. & 1 semaine \\
\bottomrule
\end{tabular}
\end{table}

\section*{Conclusion}
\addcontentsline{toc}{section}{Conclusion}
L'étude de l'existant a confirmé l'intérêt d'une solution intégrée, accessible et adaptée aux petites et moyennes structures. Les objectifs et le périmètre définis dans ce chapitre délimitent précisément ce que Tastify doit couvrir. Le chapitre suivant traduit ces besoins en une conception fonctionnelle et objet à l'aide d'UML.
